\section{Mathematical Resolution and Computational Bounds}

Establishing topological invariance is a prerequisite; the subsequent requirement is demonstrating that the generated braid group representations are dense in the special unitary group, thereby achieving universal quantum computation capabilities.

\subsection{Density in the Special Unitary Group and Solovay-Kitaev}
The unitary representations of the braid group $B_n$ generated by the transformations among the 24 Atlas classes map into $SU(2^n)$. To claim universality, the generated set of unitaries must be capable of approximating any target unitary transformation. The conceptual model dictates the specific algebraic generators derived from the symmetry classes. 

Because the eigenspaces of these generators involve phases that are incommensurate with rational multiples of $\pi$, the group they generate is dense in $SU(2)$. According to the Solovay-Kitaev theorem, this density guarantees that any arbitrary quantum gate can be approximated to precision $\epsilon$ using a sequence of operations of length $O(\log^c(1/\epsilon))$. This mathematical resolution is explicitly validated by the repository's test matrix, which constructs sample approximation sequences to bounded precision against external F1 mathematical oracles.

\subsection{Delineation of Complexity Bounds}
It is crucial to delineate the complexity-theoretic boundaries of this implementation. The UOR Atlas executes deterministically on classical computational hardware (the holospace substrate). The classical simulation of generic quantum circuits is bounded by exponential overhead in time or space; therefore, the UOR Atlas does not claim to execute arbitrary polynomial-time quantum algorithms (the BQP complexity class) efficiently in polynomial classical time (P), as such a claim would imply $P = BQP$.

Instead, the framework acts as an exact, algebraically invariant validation ledger and topological compiler. It allows researchers to formally design, compile, and algebraically verify universal topological algorithms without the noise interference inherent to physical qubits. The execution of these algorithms scales classically, bounded by the tensor network contraction limits of the host machine, while preserving absolute logical fidelity.
